% Figure 1: what the study does, in the order of the paper. Every number is a macro produced by the pipeline, so the
% figure cannot drift from the results; step 11 verifies that each one exists.
\newcommand{\boxheight}{4.75cm}
\begin{figure}[!h]
	\centering
	\resizebox{\textwidth}{!}{%
		\begin{tikzpicture}[
			font=\sffamily\scriptsize,
			>=Stealth,
			node distance=0.45cm,
			basebox/.style={
				thick,
				rounded corners=4pt,
				inner sep=6pt,
				anchor=center,
				align=center
			},
			phase1/.style={basebox, draw=red!70!black, fill=red!4},
			phase2/.style={basebox, draw=blue!75!black, fill=blue!4},
			phase3/.style={basebox, draw=teal!80!black, fill=teal!4},
			phase4/.style={basebox, draw=violet!75!black, fill=violet!4},
			phase5/.style={basebox, draw=gray!70!black, fill=gray!5},
			arrowline/.style={->, ultra thick, draw=gray!65}
			]

			% STAGE 1: AUDIT OF THE PUBLISHED PIPELINE
			\node[phase1] (box1) {
				\begin{minipage}[t][\boxheight][t]{3.15cm}
					\centering
					\textbf{\footnotesize 1. Audit} \par\vspace{2pt}\hrule\vspace{7pt}
					\textbf{Object:} \\published code, re-run as released \par\vspace{4pt}
					\textbf{Leak classes:} \\re-split after augmentation; early stopping on test; in-sample importance; duplicates \par\vspace{4pt}
					\textbf{Removed one at a time:} \\$R^2$ \KarLadderRungZeroRTwoMean\ $\rightarrow$ \KarLadderRungFourRTwoMean
				\end{minipage}
			};

			% STAGE 2: LEAKAGE-FREE BENCHMARK
			\node[phase2, right=of box1] (box2) {
				\begin{minipage}[t][\boxheight][t]{3.15cm}
					\centering
					\textbf{\footnotesize 2. Benchmark} \par\vspace{2pt}\hrule\vspace{7pt}
					\textbf{Data:} \\$N = \NCompounds$ TgDHFR \par\vspace{4pt}
					\textbf{Protocol:} \\$\OuterKFolds \times \OuterRepeats$ nested CV, \TotalOuterFolds\ shared folds \par\vspace{4pt}
					\textbf{Models:} \\4 representations vs.\ 7 neural \par\vspace{4pt}
					\textbf{Inference:} \\Nadeau--Bengio, Holm, TOST
				\end{minipage}
			};

			% STAGE 3: ROBUSTNESS
			\node[phase3, right=of box2] (box3) {
				\begin{minipage}[t][\boxheight][t]{3.15cm}
					\centering
					\textbf{\footnotesize 3. Robustness} \par\vspace{2pt}\hrule\vspace{7pt}
					\textbf{Chance correlation:} \\y-randomization, 4 forms \par\vspace{4pt}
					\textbf{Chemotype shift:} \\scaffold-grouped splits (\ScaffoldNGroups\ groups) \par\vspace{4pt}
					\textbf{Curation:} \\extreme-potency records retained \par\vspace{4pt}
					\textbf{Cost:} \\wall-clock time per model
				\end{minipage}
			};

			% STAGE 4: FINDINGS
			\node[phase4, right=of box3] (box4) {
				\begin{minipage}[t][\boxheight][t]{3.15cm}
					\centering
					\textbf{\footnotesize 4. Findings} \par\vspace{2pt}\hrule\vspace{7pt}
					\textbf{Morgan ensemble:} \\$R^2 = \MorganRTwoMean$ \par\vspace{4pt}
					\textbf{Default D-MPNN:} \\$R^2 = \ChemPropDMPNNFoldRTwoMean$ \par\vspace{4pt}
					\textbf{Pretrained D-MPNN:} \\$R^2 = \ChemeleonRTwoMean$, the only equivalent model \par\vspace{4pt}
					\textbf{Reported gain:} \\accounted for by the leaks
				\end{minipage}
			};

			% STAGE 5: APPLICATION
			\node[phase5, right=of box4] (box5) {
				\begin{minipage}[t][\boxheight][t]{3.15cm}
					\centering
					\textbf{\footnotesize 5. Application} \par\vspace{2pt}\hrule\vspace{7pt}
					\textbf{Library:} \\$N = \ScreenNLibrary$ FDA-approved \par\vspace{4pt}
					\textbf{Applicability domain:} \\\ScreenNDomain\ of \ScreenNStructural\ inside \par\vspace{4pt}
					\textbf{Orthogonal check:} \\docking, redock RMSD \RedockRMSD\,\AA \par\vspace{4pt}
					\textit{Hypotheses, not leads; no selectivity assay}
				\end{minipage}
			};

			\draw[arrowline] (box1) -- (box2);
			\draw[arrowline] (box2) -- (box3);
			\draw[arrowline] (box3) -- (box4);
			\draw[arrowline] (box4) -- (box5);

		\end{tikzpicture}%
	}
	\caption{\textbf{Structure of the study.} The published pipeline is first re-executed and taken apart one leak at a
		time, which measures what each class of leak contributes to its reported $R^2$ (1). The same data are then used
		for a leakage-free benchmark in which every model is scored on identical outer folds and every comparison is
		corrected for the overlap of training sets and tested for equivalence (2), and the conclusions are checked
		against chance correlation, a harder splitting scheme, the curation decision and computational cost (3). The
		findings (4) are that a fingerprint ensemble is not outperformed by a graph network at this data volume, that
		only a pretrained network reaches equivalence with it, and that the gain reported by the original study is
		accounted for by the leaks removed in (1). The final model is applied to approved drugs as hypothesis
		generation, with each candidate placed relative to the applicability domain (5).}
	\label{fig:pipeline_flowchart}
	\vspace{-10pt}
\end{figure}
